% !TEX program = lualatex
% Transparencias de Fundamentos de las Matemáticas II
% Clase 06: Ecuaciones diofánticas lineales

\documentclass[aspectratio=169,11pt]{beamer}

\usepackage{fontspec}
\usepackage[spanish,es-nodecimaldot]{babel}
\usepackage{amsmath,amssymb,mathtools}
\usepackage{booktabs}
\usepackage{csquotes}
\usepackage{xcolor}

\usetheme{Madrid}
\usecolortheme{default}
\setbeamertemplate{navigation symbols}{}
\setbeamertemplate{footline}[frame number]

\definecolor{FdMBlue}{RGB}{20,55,100}
\definecolor{FdMRed}{RGB}{130,30,30}
\definecolor{FdMGreen}{RGB}{25,100,60}

\setbeamercolor{structure}{fg=FdMBlue}
\setbeamercolor{title}{fg=white,bg=FdMBlue}
\setbeamercolor{frametitle}{fg=white,bg=FdMBlue}
\setbeamercolor{block title}{fg=white,bg=FdMBlue}
\setbeamercolor{block body}{bg=blue!3}
\setbeamercolor{alerted text}{fg=FdMRed}

\setmainfont{Latin Modern Roman}
\setsansfont{Latin Modern Sans}
\setmonofont{Latin Modern Mono}

\newcommand{\N}{\mathbb{N}}
\newcommand{\Z}{\mathbb{Z}}
\newcommand{\Q}{\mathbb{Q}}
\newcommand{\R}{\mathbb{R}}
\newcommand{\C}{\mathbb{C}}
\newcommand{\Pow}{\mathcal{P}}
\newcommand{\id}{\operatorname{id}}
\newcommand{\mcd}{\operatorname{mcd}}
\newcommand{\mcm}{\operatorname{mcm}}
\newcommand{\im}{\operatorname{Im}}

\title[FdM II -- Clase 06]{Ecuaciones diofánticas lineales}
\subtitle{Soluciones enteras de \(ax+by=c\)}
\author{Antonio Falcó Montesinos}
\institute{Universidad CEU Cardenal Herrera}
\date{Fundamentos de las Matemáticas II}

\begin{document}

\begin{frame}
  \titlepage
\end{frame}

\begin{frame}{Objetivos de la clase}
\begin{itemize}
  \item Formular ecuaciones diofánticas lineales.
  \item Usar el criterio de existencia de soluciones.
  \item Obtener soluciones a partir de Bézout.
\end{itemize}
\end{frame}

\begin{frame}{Mapa de la clase}
\tableofcontents
\end{frame}

\section{Planteamiento}

\begin{frame}{Planteamiento}
\begin{block}{Ecuación diofántica}
\begin{itemize}
  \item Una ecuación diofántica lineal tiene la forma \(ax+by=c\).
  \item Los coeficientes \(a,b,c\) son enteros.
  \item Se buscan soluciones \((x,y)\in\mathbb{Z}^2\).
\end{itemize}
\end{block}
\end{frame}

\begin{frame}{Planteamiento}
\begin{block}{Criterio}
\begin{itemize}
  \item La ecuación \(ax+by=c\) tiene solución entera si y solo si \(\mcd(a,b)\mid c\).
  \item Si el \(\mcd\) no divide a \(c\), no hay solución.
  \item Si divide a \(c\), Bézout permite construir una.
\end{itemize}
\end{block}
\end{frame}

\begin{frame}{Planteamiento}
\begin{block}{Ejemplo negativo}
\begin{itemize}
  \item \(6x+10y=7\) no tiene solución entera.
  \item \(\mcd(6,10)=2\).
  \item Como \(2\nmid 7\), el criterio descarta soluciones.
\end{itemize}
\end{block}
\end{frame}

\section{Construcción}

\begin{frame}{Construcción}
\begin{block}{Ejemplo positivo}
\begin{itemize}
  \item Resolver \(18x+30y=6\).
  \item Como \(\mcd(18,30)=6\), hay solución.
  \item Una combinación es \(6=30-18\), luego \(x=-1\), \(y=1\).
\end{itemize}
\end{block}
\end{frame}

\begin{frame}{Construcción}
\begin{block}{Familia de soluciones}
\begin{itemize}
  \item Una solución particular no suele ser única.
  \item A partir de una solución se obtiene una familia infinita.
  \item El estudio completo requiere parametrizar usando \(a/d\) y \(b/d\).
\end{itemize}
\end{block}
\end{frame}

\section{Profundización}

\begin{frame}{Profundización}
\begin{block}{Lectura estructural}
\begin{itemize}
  \item Bézout resuelve el caso \(c=d\).
  \item Multiplicar la combinación por \(c/d\) resuelve el caso general.
  \item Esto explica el papel exacto de la divisibilidad \(d\mid c\).
\end{itemize}
\end{block}
\end{frame}

\begin{frame}{Profundización}
\begin{block}{Errores frecuentes}
\begin{itemize}
  \item Buscar soluciones por tanteo sin comprobar el criterio.
  \item Confundir solución racional con solución entera.
  \item Olvidar que el criterio usa \(\mcd(a,b)\).
\end{itemize}
\end{block}
\end{frame}

\begin{frame}{Profundización}
\begin{block}{Pregunta de control}
\begin{itemize}
  \item ¿Tiene solución \(14x+21y=35\)?
  \item ¿Y \(14x+21y=5\)?
  \item Justificar cada respuesta antes de buscar valores.
\end{itemize}
\end{block}
\end{frame}


\section{Detalle y práctica}

\begin{frame}{Detalle y práctica}
\begin{block}{Familia completa de soluciones}
\begin{itemize}
  \item Si \(d=\mcd(a,b)\) y \(d\mid c\), una solución particular \((x_0,y_0)\) genera todas las soluciones.
  \item La forma general es \(x=x_0+\frac{b}{d}t\), \(y=y_0-\frac{a}{d}t\), con \(t\in\mathbb{Z}\).
  \item El parámetro \(t\) recoge la libertad que queda tras imponer una sola ecuación.
\end{itemize}
\end{block}
\end{frame}

\begin{frame}{Detalle y práctica}
\begin{block}{Actividad breve}
\begin{itemize}
  \item Resolver \(12x+18y=30\).
  \item Encontrar una solución particular usando Bézout.
  \item Escribir la familia completa y comprobar dos valores distintos de \(t\).
\end{itemize}
\end{block}
\end{frame}

\begin{frame}{Cierre}
\begin{itemize}
  \item El criterio \(\mcd(a,b)\mid c\) decide existencia.
  \item Bézout construye soluciones.
  \item Las ecuaciones diofánticas conectan divisibilidad y álgebra.
\end{itemize}
\end{frame}

\begin{frame}{Trabajo recomendado}
\begin{itemize}
  \item Releer las secciones correspondientes del manual.
  \item Identificar definiciones, símbolos nuevos y enunciados principales.
  \item Preparar dudas y ejemplos para el seminario asociado.
\end{itemize}
\end{frame}

\end{document}
